\rtmtight
\begin{tblr}{width=\textwidth,colspec={|Q[c,m,2.1cm]|X[l,m]|},hlines,vlines,hspan=minimal,rowsep=1.5pt,colsep=3pt,row{1}={bg=headgray,font=\bfseries}}
\SetCell{c,m}\hfil\logo\hfil & \textbf{POLITEKNIK NEGERI MALANG}\par\textbf{JURUSAN TEKNOLOGI INFORMASI}\par\textbf{PROGRAM STUDI: D4 TEKNIK INFORMATIKA} \\
\SetCell[c=2]{l} \textbf{RENCANA TUGAS MAHASISWA (RTM) DAN RUBRIK PENILAIAN} & \\
Nama Mata Kuliah & Big Data \\
Kode Mata Kuliah & RTI256106 \quad \textbf{sks / Jam perminggu:} 2 SKS/4 jam \quad \textbf{Semester:} 6 \\
Dosen Pengampu & Tim Pengajar Program Studi D4 TI \\
\end{tblr}
\assessmentblock{Tugas Pipeline Big Data}{Case Method dan praktik Big Data}{SCPMK0702-05001, SCPMK0904-05002}{Mahasiswa merancang dan mengimplementasikan pipeline pemrosesan Big Data serta menganalisis arsitektur sistem terdistribusi berdasarkan studi kasus yang diberikan.}{Menganalisis kebutuhan data, merancang arsitektur pipeline, mengimplementasikan menggunakan Hadoop/Spark, menguji throughput, dan mendokumentasikan.}{Kode pipeline (Hadoop/Spark), dokumen arsitektur, laporan analisis performa, dan bahan presentasi.}{Ketepatan rancangan arsitektur pipeline Big Data & 12 & Arsitektur pipeline mencakup ingestion, processing, storage, dan serving layer dengan justifikasi teknologi yang tepat dan sesuai kebutuhan skala data. & Arsitektur tersedia dan mencakup komponen utama tetapi beberapa layer atau justifikasi teknologi masih kurang detail. & Arsitektur tidak mencakup semua layer pipeline atau tidak dapat dijadikan acuan implementasi. \\ \\
Implementasi pipeline menggunakan Hadoop/Spark & 12 & Pipeline diimplementasikan dengan benar, job berjalan tanpa error, data terproses sesuai spesifikasi, dan output konsisten. & Pipeline sebagian berjalan tetapi beberapa job mengalami error, data tidak konsisten, atau performa di bawah ekspektasi. & Pipeline tidak dapat berjalan atau output tidak sesuai dengan spesifikasi yang diminta. \\ \\
Analisis performa dan dokumentasi & 6 & Laporan mencakup analisis throughput, latency, resource utilization, dan rekomendasi optimasi yang konkret. & Laporan tersedia tetapi analisis performa atau rekomendasi optimasi masih kurang mendalam. & Tidak ada analisis performa yang bermakna atau dokumentasi tidak mencerminkan implementasi yang dilakukan. \\}

\assessmentblock{Kuis Konsep dan Arsitektur Big Data}{Kuis}{SCPMK0904-05002, SCPMK0702-05001}{Kuis tertulis untuk mengukur pemahaman konsep Big Data, arsitektur sistem terdistribusi, dan teknologi pemrosesan data skala besar.}{Menjawab soal pilihan ganda dan/atau uraian terkait konsep 4V, arsitektur Lambda/Kappa, HDFS, MapReduce, Spark, dan streaming.}{Jawaban tertulis kuis.}{Pemahaman konsep dan arsitektur Big Data & 5 & Konsep 4V, arsitektur Lambda/Kappa, HDFS, dan prinsip distribusi dijelaskan dengan benar dan akurat. & Sebagian besar konsep dipahami tetapi penjelasan arsitektur atau prinsip distribusi masih kurang tepat. & Konsep Big Data tidak dipahami atau jawaban tidak relevan dengan materi yang diujikan. \\ \\
Penguasaan teknologi pemrosesan Big Data & 5 & MapReduce, Spark RDD/DataFrame, dan Kafka dijelaskan dengan tepat termasuk perbedaan batch dan stream processing. & Sebagian teknologi dipahami tetapi perbandingan atau penjelasan mekanisme kerja masih kurang akurat. & Teknologi pemrosesan Big Data tidak dapat dijelaskan atau pemahaman fundamental keliru. \\}

\assessmentblock{UTS Arsitektur dan MapReduce}{UTS berbasis analisis dan implementasi}{SCPMK0904-05002, SCPMK0702-05001}{Evaluasi tengah semester untuk mengukur kemampuan menganalisis arsitektur Big Data dan mengimplementasikan pemrosesan data menggunakan paradigma MapReduce atau Spark dasar.}{Menganalisis studi kasus arsitektur Big Data, mengimplementasikan program MapReduce atau Spark job dasar, dan menjelaskan keputusan teknis.}{Kode MapReduce/Spark job, dokumen analisis arsitektur, dan penjelasan keputusan teknis.}{Kemampuan analisis arsitektur Big Data & 8 & Arsitektur yang tepat dipilih dan dianalisis dengan menjelaskan trade-off Lambda vs Kappa, pertimbangan skalabilitas, dan kesesuaian dengan kasus. & Analisis arsitektur tersedia tetapi trade-off atau justifikasi pemilihan masih kurang komprehensif. & Arsitektur tidak dapat dianalisis dengan benar atau pemilihan tidak sesuai dengan kebutuhan kasus. \\ \\
Implementasi MapReduce atau Spark job & 7 & Program MapReduce/Spark diimplementasikan dengan benar, logika map/reduce/transform tepat, dan output sesuai dengan yang diharapkan. & Implementasi sebagian besar benar tetapi logika map/reduce atau transformasi data masih ada kesalahan. & Program tidak berjalan atau logika fundamental MapReduce/Spark tidak diimplementasikan dengan benar. \\}

\assessmentblock{PBL Implementasi Pipeline Big Data}{Project Based Learning}{SCPMK0702-05001, SCPMK0503-05003}{Mahasiswa mengimplementasikan pipeline pemrosesan Big Data lengkap menggunakan ekosistem Hadoop/Spark yang mencakup batch processing, stream processing, dan penyimpanan data dalam data lakehouse, sebagai solusi komputasi paralel dan terdistribusi berbasis cloud.}{Membangun pipeline end-to-end, mengimplementasikan Spark job untuk batch dan streaming, mengintegrasikan Kafka, menyimpan ke Delta Lake/Iceberg, dan memvisualisasikan hasil.}{Pipeline berjalan di cluster/Docker, kode sumber, data catalog, laporan implementasi, dan visualisasi output data.}{Kelengkapan dan fungsionalitas pipeline & 7 & Pipeline mencakup ingestion, batch processing, stream processing, storage, dan serving; semua komponen berjalan end-to-end tanpa error kritis. & Pipeline mencakup sebagian besar komponen tetapi beberapa stage (terutama streaming atau storage) belum berjalan optimal. & Pipeline tidak dapat berjalan end-to-end atau lebih dari setengah komponen tidak berfungsi. \\ \\
Kualitas implementasi Spark dan Kafka & 7 & Spark job teroptimasi (partitioning, caching, serialization), Kafka consumer/producer terkonfigurasi dengan benar, dan throughput memenuhi target. & Spark job dan Kafka berjalan tetapi optimasi atau konfigurasi masih dapat ditingkatkan secara signifikan. & Spark job atau Kafka tidak berjalan dengan benar atau throughput sangat jauh di bawah target. \\ \\
Dokumentasi dan visualisasi output & 3 & Kode terdokumentasi, arsitektur pipeline divisualisasikan, output data divisualisasikan secara bermakna, dan laporan implementasi lengkap. & Dokumentasi dan visualisasi tersedia tetapi beberapa komponen masih kurang rinci atau tidak informatif. & Dokumentasi tidak mencukupi untuk mereproduksi pipeline atau output data tidak divisualisasikan. \\ \\
Solusi komputasi paralel dan terdistribusi berbasis cloud & 5 & Solusi komputasi paralel-terdistribusi dirancang dan diterapkan secara tepat untuk kebutuhan Big Data dan platform cloud. & Solusi tersedia, tetapi pemanfaatan sumber daya paralel-terdistribusi atau cloud belum optimal. & Solusi tidak memanfaatkan komputasi paralel-terdistribusi atau tidak sesuai kebutuhan cloud. \\}

\assessmentblock{UAS Demo dan Analisis Pipeline}{UAS berbasis demo dan presentasi}{SCPMK0702-05001, SCPMK0904-05002}{Mahasiswa mendemokan pipeline Big Data yang telah diimplementasikan, menganalisis arsitektur dan performa, serta mempertahankan keputusan teknis secara komprehensif.}{Demo pipeline live, presentasi arsitektur dan analisis performa, menjawab pertanyaan teknis dari penilai.}{Demo pipeline live, laporan analisis performa, dan bahan presentasi.}{Kualitas demo pipeline Big Data & 9 & Demo menampilkan pipeline end-to-end berjalan dengan data nyata, throughput terukur, dan output divisualisasikan secara live. & Demo sebagian besar berjalan tetapi beberapa komponen tidak dapat didemonstrasikan atau data tidak akurat. & Demo tidak dapat menampilkan pipeline yang berjalan atau output tidak relevan. \\ \\
Analisis arsitektur dan performa pipeline & 8 & Analisis mencakup throughput, latency, skalabilitas, bottleneck, dan rekomendasi optimasi berbasis data pengukuran aktual. & Analisis tersedia tetapi beberapa metrik, identifikasi bottleneck, atau rekomendasi optimasi masih kurang mendalam. & Analisis tidak berdasarkan data aktual atau tidak dapat mengidentifikasi bottleneck dan area optimasi. \\ \\
Defense teknis dan pemahaman arsitektur & 6 & Keputusan arsitektur dan teknologi dijelaskan dengan justifikasi teknis yang kuat, trade-off dipahami, dan pertanyaan dijawab kompeten. & Defense cukup baik tetapi beberapa keputusan arsitektur atau trade-off belum dapat dipertahankan dengan baik. & Tidak dapat menjelaskan keputusan arsitektur atau menjawab pertanyaan teknis fundamental Big Data. \\}

\begin{tblr}{width=\textwidth,colspec={|Q[l,m,3.2cm]|X[l,m]|},hlines,vlines,hspan=minimal,row{1-3}={bg=headgray},rowsep=1.5pt,colsep=3pt}
Jadwal Pelaksanaan: & Minggu 1--17 sesuai RPS. Setiap evaluasi memiliki RTM dan rubrik multi-kriteria. \\
Lain-Lain Yang Diperlukan: & LMS, repositori/kode sumber pipeline, perangkat lunak sesuai mata kuliah, dan perangkat presentasi. \\
Pustaka: & Mengacu pustaka utama dan pendukung pada bagian RPS. \\
\end{tblr}
